% Capítulo incluído por tcc-ufv.tex com \include: sem preâmbulo e sem \begin{document}.
\chapter{Referencial teórico}\label{cap:referencial}

Segundo \textcite{ex-livro}, a citação direta longa (mais de três linhas)
fica recuada:
\begin{citacao}
Texto citado literalmente, com mais de três linhas, que a classe formata
com recuo de 4\,cm, fonte menor e espaçamento simples, como orienta a NBR
10520. A referência vem ao final da citação, entre parênteses
\cite[p.~35]{ex-livro}.
\end{citacao}

Uma \ac{api} define como dois sistemas trocam dados; as \acp{api} da
\autoref{fig:arquitetura} usam arquivos \ac{csv} \cite{ex-artigo}.

% ilustracao: mede a imagem e alinha título, fonte, legenda e nota na largura dela
\begin{figure}[htbp]
\begin{ilustracao}{\caption{Arquitetura do protótipo}\label{fig:arquitetura}}%
                  {\elaboracaopropria
                   \legenda{A letra A representa o servidor.}}
  \includegraphics[width=.6\linewidth]{example-image-a}
\end{ilustracao}
\end{figure}
